% SPDX-License-Identifier: LPPL-1.3c
% Copyright (C) 2026 Christophe Jorssen
% Author and Current Maintainer: Christophe Jorssen <christophe.jorssen@gmail.com>
% LPPL maintenance status: maintained
% This file is part of luafemm. See LICENSE and LICENSES.md for its terms.
\section{Ideal magnetic circuits and mean contours}
\label{sec:ideal}
An ideal circuit confines magnetic flux to a prescribed union of iron,
permanent magnets and air gaps. No flux leaves the lateral walls, including
the lateral walls of the gaps. This is a different physical boundary-value
problem from the ordinary open-air model: clipping the ordinary field picture
would not impose this condition.

Flux is conserved along every unbranched tube. At a junction the signed
incoming and outgoing fluxes balance. Neither statement requires a uniform
local field. A narrowing section increases its average normal induction,
and the local field also varies around corners and inside curved material.
The material laws remain finite and may be nonlinear. The word \emph{ideal}
does not replace iron by infinite permeability or remove its magnetomotive
force drop.

\subsection{A first confined circuit}
Let us begin with a complete U electromagnet. The node already knows where
its iron, winding sections and two gaps are located. In ideal mode it includes
the gaps in the computational domain, even if |gap mesh size=0|. It excludes
the winding sections and uses their enclosed currents as excitation.

\begin{codeexample}[width=6cm]
\begin{tikzpicture}[
  femm/problem={field model=ideal,
    depth=10,mesh size=4},scale=.55]
  \node[femm/u electromagnet={
    ampere turns=1200,
    gap mesh size=2}] (M) {};
  \pic[femm/lines=9] {femm field};
  \pic[femm/mean={component=M}]
    {femm mean path};
\end{tikzpicture}
\end{codeexample}

The red contour is a geometric reference for applying Ampere's law. It follows
the centres of the component sections and closes through the gaps. It need
not coincide with one of the blue field lines. Start with a coarse mesh while
arranging the figure; refine the mesh and curved geometry when examining local
field values or preparing the final result. Conservation of flux is a discrete
identity, but it does not by itself measure the accuracy of the predicted flux.

\begin{key}{/femm/problem/field model=\meta{open or ideal} (initially open)}
Choose the ordinary exterior-air model or the confined-domain model through
PGF's choice handler. Set it inside |femm/problem| before the picture begins.
|ideal| requires |mesher=delaunay| and at least one included physical region.
The bounding rectangle still encloses every declared vertex and bounds the
staging mesh, but its exterior air is removed before assembly.
Default rectangular boundary data are unused; nondefault exterior data,
Lua source/boundary callbacks and FEM import are rejected in this mode.
The option follows ordinary TeX scoping and participates in cache identity.
\end{key}

\begin{key}{/femm/region/ideal domain=\meta{true or false} (initially false)}
Mark a region as included in the confined computational domain. Supply it
inside |femm/region|; it inherits within the surrounding TeX scope. The last
covering region determines both the material and this membership flag.
Thus an air gap or a later material overlay needs |ideal domain=true| as well.
Compound-path holes preserve earlier declarations, as for ordinary materials.
The flag has no effect in |field model=open|.

Component nodes set the flag for their core, armature, magnet and gap parts,
and clear it for winding parts. A current density inside the active ideal
domain is rejected: windings must lie in excluded windows or exterior air.
All physical paths still accept the normal TikZ curves and transformations.
\end{key}

For a native path model, give the iron and the confined gaps their own styles.
The following complete example uses a rounded compound path, an explicit
window excitation and a user-drawn mean contour. Its gap section is also a
named profile, allowing the flux to be printed without numerical quadrature.
\luafemmexample{ideal-paths}

\subsection{Excitation and declaration order}
\begin{key}{/tikz/femm/excitation=\marg{options}}
Add a signed enclosed-current seed for ideal mode. Each declaration resets its
|x|, |y| and |ampere turns| options to zero before processing the supplied list.
Picture options and outer-scope styles queue their declarations until the
problem exists, independently of their position relative to |femm/problem|.
Inside a started picture the declaration is immediate. Several declarations
add; they do not replace one another. The seed is passive data, not a Lua
expression or a request to draw a winding.
\end{key}
\begin{keylist}[/femm/excitation]{x=\meta{number},y=\meta{number}}
Seed coordinates in model millimetres, initially zero for each declaration.
These are numerical model coordinates, unaffected by a local drawing transform.
At meshing time the point must lie strictly inside an excluded bounded window;
points on interfaces, inside magnetic material or in unbounded exterior air
are errors.
\end{keylist}
\begin{key}{/femm/excitation/ampere turns=\meta{number} (initially 0)}
Signed enclosed $NI$ in amperes. Positive values point along $+z$, and produce
positive counterclockwise circulation of $\mathbf H$ around that window.
Negative and zero values are accepted. Seeds add to currents already described
by winding regions. Do not add a seed for a winding that a component has
already declared.
\end{key}

The order is: select the problem and cache policy; declare all components,
physical paths and excitations; then request a field, a mean contour, a mesh,
a profile plot or a flux value. Every one of those requests can freeze the
physical model. Merely registering a |femm/profile| path still does not solve.
Named profiles retain their originating model after the picture ends.
Generic component nodes have exactly the same placement rules in both modes.

\subsection{Mean contours and transported-flux lines}
\begin{key}{/tikz/pics/femm mean path}
Resolve the component and cycle selected by |femm/mean|, solve on demand and
draw the resulting contour in the model's original picture frame. Local pic
translations do not detach it from the component. An unknown component,
unavailable cycle or contour outside the computational domain is an error.
\end{key}
\begin{key}{/tikz/femm/mean=\marg{options}}
Set the following options in the |/femm/mean| family. They inherit according
to normal PGF scoping; a pic's local options do not leak into the next pic.
\end{key}
\begin{key}{/femm/mean/component=\meta{node name} (initially empty)}
Required name of an already declared component node in the current model.
Use the full node name when applying a TikZ name prefix or suffix.
\end{key}
\begin{key}{/femm/mean/cycle=\meta{name} (initially main)}
Select |main| on a U electromagnet or either toroid; select |left| or |right|
on an E electromagnet. Bare U and E cores have no closed mean cycle.
The two E cycles share the central branch. They are contours for Ampere's law;
their shared segment is not a bifurcating field line.
\end{key}
\begin{key}{/femm/mean/kind=\meta{geometric or flux} (initially geometric)}
|geometric| follows section centres, orienting the contour with its net
circulation of $\mathbf B$. At zero field it keeps its construction direction.
Changing the excitation sign can therefore reverse its profile's starting
direction. No implication of uniform $\mathbf H$ along this contour is made.

|flux| selects an actual computed isoline halfway between the selected window
potential and its outer reference potential. It requires ideal mode, one
enclosed window and one unambiguous closed isoline surrounding that window
alone. It rejects a zero-flux window or a branched situation without such an
isolated line; it does not silently substitute a geometric contour. In a
single tube this line divides the transported flux equally, rather than
necessarily dividing the geometric section equally. Its starting point follows
the contour extractor's mesh ordering and can change after remeshing.
\end{key}
\begin{key}{/femm/mean/profile=\meta{identifier} (initially empty)}
Optionally register the chosen contour for |\femmplot|. Empty means no profile.
The identifier follows the existing profile rules; reusing it replaces the
previous profile. The normal offset is zero and outside-domain samples error.
\end{key}
\begin{key}{/femm/mean/samples=\meta{integer} (initially 201)}
Number of equally spaced arc-length profile samples, from 2 to 100000.
It changes only sampling, not the field mesh or the shape of the drawn contour.
\end{key}
\begin{stylekey}{/tikz/femm mean path}
The base style uses |draw=red!75!black|, |thick|, |dashed| and
|femm oriented field|. Use |/.append style| to retain the direction markers. The existing
arrow spacing, dimensions and enable/disable keys apply.
\end{stylekey}
\begin{key}{/tikz/femm/mean style=\marg{TikZ options}}
Set additional options applied after the base mean style, initially empty.
This replaces the extra options locally. |femm current mean| is internal
storage, not a physical declaration.
\end{key}

For a custom contour, use an ordinary closed |\draw| path with the
|femm mean path| style and |femm/profile={name=...}|. This keeps arbitrary
Bezier curves, rounded corners and named coordinates available. Its direction
is the direction in which the user writes the path; it is not automatically
reoriented. Automatic skeleton extraction from an arbitrary material domain
is deliberately not assumed.

\subsection{Flux values and Ampere profiles}
\begin{command}{\femmflux\oarg{options}}
Solve the source model of a named profile and print its signed section flux
in webers using PGF's standard number printer. Options belong to |/femm/flux|.
The command works in all three formats and requires a positive problem |depth|.
PGF's |/pgf/number format| styles control numerical presentation.
\end{command}
\begin{key}{/femm/flux/profile=\meta{identifier}}
Required named section profile; empty at the start of each command invocation.
Its path must remain in the computational domain and have |offset=0|. Declare
a displaced section explicitly if needed. For a path from $P$ to $Q$, with
right normal, the returned value is
\[
 \Phi_{P\to Q}=d\,[A_z(Q)-A_z(P)],
\]
where $d$ is the physical extrusion depth in metres. Reversing the path reverses
the sign. This right-normal convention is the opposite of the left normal
used by the profile component |Bn|. The result depends on endpoint potentials,
not on the number of samples. A closed section path returns zero.
\end{key}

\begin{key}{/femm/flux/snap tolerance=\meta{nonnegative number} (initially .003)}
Endpoint projection tolerance in model millimetres. Before integrating, project
each endpoint to the nearest computational boundary if it is closer than this
distance. This accommodates small PGF coordinate differences between rotated
node anchors and captured part paths. Intermediate vertices and the physical
mesh are unchanged; the adjusted section must still remain in the domain.
Use zero to disable projection, or a smaller tolerance for very thin features.
The Lua |flux| API defaults to zero and accepts this tolerance as its third
argument. Like other option values, the TeX tolerance inherits in its scope.
\end{key}

To compare branch fluxes, place each section all the way between its two walls
and use consistent orientations. Use actual component pole anchors or captured
path endpoints: curved boundaries are polygonalised, and a point on the ideal
analytic circle can lie just outside the computational polygon.

\begin{key}{/femm/plot/component=circulation}
Additional PGFPlots component, in amperes:
$C(s)=\int_0^s\mathbf H\cdot d\boldsymbol\ell$ along the named profile.
The integral is evaluated by splitting each polyline segment at triangle
edges, then integrating the piecewise-constant element field exactly.
Increasing |samples| adds plotted abscissas; it does not improve the integral.
An offset profile or a path leaving the domain is rejected for this component,
even with |outside=nan|. At an exact material/element interface the usual first
containing triangle convention applies.

A closed mean contour should approach the linked $NI$ as the mesh is refined.
Finite-element $\mathbf H$ is not exactly tangentially continuous, so an
arbitrary contour's circulation need not equal $NI$ exactly on a coarse mesh.
This is distinct from the exact discrete conservation of magnetic flux.
\end{key}

Here are complete U examples with constant and unequal sections. Both show
the field, the geometric contour and its profile; the last plot displays
the magnetomotive force accumulated along the unequal-section circuit.
\luafemmexample{ideal-u}

\subsection{Additional component shapes}
All keys below extend |/femm/component|. They use numeric model millimetres
unless described as angles. They inherit within the surrounding TeX scope;
options on a node are local to that node. The existing material selectors,
mesh sizes, painting styles, affine placement and fixed node envelopes apply.
The E and toroidal shapes work in both field models. In open mode an air-gap
refinement region is optional; in ideal mode it is always included.

\begin{keylist}[/tikz/femm]{e core=\marg{options},e electromagnet=\marg{options}}
Select an E core open towards local $+y$, with a central leg at $x=0$.
The electromagnet adds an armature, two winding sections around the central
leg and up to three gaps. Positive |ampere turns| puts positive $J_z$ in
the left winding section and negative $J_z$ in the right. The core alone
adds no winding, armature, excitation or mean cycle.
\end{keylist}
\begin{keylist}[/femm/component]{left thickness=\meta{number},
 right thickness=\meta{number},yoke thickness=\meta{number},
 center thickness=\meta{number}}
Initially zero: inherit |leg thickness|. Positive overrides specify the
individual section thicknesses. Left, right and yoke apply to U and E shapes;
centre applies only to E. Negative values are errors. The resulting windows
and free leg heights must remain positive, and the winding must fit.
Changing one section does not change the independently specified armature.
\end{keylist}
\begin{keylist}[/femm/component]{left gap=\meta{number},
 center gap=\meta{number},right gap=\meta{number}}
E-only gap lengths, initially $-1$ to inherit |gap|. Otherwise require
nonnegative values; zero joins that pole directly to the armature.
The armature underside is at |height/2+max(left gap,center gap,right gap)|,
after resolving inherited values. Each leg stops its own gap distance below
this plane. A gap cannot consume the full leg above the yoke. These options
do not change the equal-gap U component.
\end{keylist}

The E shape provides the usual envelope and |core center| anchors, plus
|left pole|, |center pole|, |right pole|, |left window|, |right window|,
|yoke center| and |window center| (the left window's text position).
The outer legs also provide |left pole inner|, |left pole outer|,
|right pole inner| and |right pole outer|. The electromagnet additionally
provides |left gap|, |center gap|, |right gap|, the corresponding three
|... armature| anchors, |armature center| and both |coil positive| and
|coil negative|. Zero-length gaps still have a centre anchor.
An unavailable anchor reports an error rather than a guessed location.

\luafemmexample{ideal-e}

\begin{keylist}[/tikz/femm]{toroid=\marg{options},tapered toroid=\marg{options}}
Select an annular core centred at |core center|, with its gap towards local
$+x$. The tapered shape adds two straight-sided transitions adjacent to the
gap. A positive |magnet span| replaces the opposite arc with permanent magnet.
These are planar cross-sections with a constant extrusion depth. The taper
can illustrate the meridional outline of truncated-conical pole pieces, but
the calculation uses planar sections, not circular areas $\pi r^2$ and not
an axisymmetric or three-dimensional field model.
\end{keylist}
\begin{key}{/femm/component/radius=\meta{number} (initially 30)}
Mean radius. Must exceed half the core thickness.
\end{key}
\begin{key}{/femm/component/core thickness=\meta{positive number} (initially 10)}
Radial width of the regular toroidal section. Its physical area is this width
times the specified constant depth; it is not a diameter of a round section.
\end{key}
\begin{key}{/femm/component/pole thickness=\meta{positive number} (initially 6)}
Width at the gap on a tapered toroid, no larger than |core thickness|.
It is unused on an ordinary toroid or when |taper angle=0|.
\end{key}
\begin{key}{/femm/component/gap angle=\meta{number} (initially 8)}
Angular opening in degrees, from zero inclusive to 90 exclusive. Pole centres
are at angles $\pm\texttt{gap angle}/2$. The gap is the quadrilateral joining
their inner and outer endpoints. Zero closes the toroid. Unlike a U gap,
this is an angle: the pole-centre separation is
$2R\sin(\texttt{gap angle}/2)$.
\end{key}
\begin{key}{/femm/component/taper angle=\meta{number} (initially 25)}
Angular extent in degrees of each transition next to a tapered toroid's gap,
from zero inclusive to 90 exclusive. Straight sides join the narrow pole
section to the full-width annulus. The geometric mean follows their centre
chord. Zero disables tapering. It is unused by an ordinary toroid.
\end{key}
\begin{key}{/femm/component/magnet span=\meta{number} (initially 0)}
Angular extent of the magnet arc, centred at $180^\circ$, from zero inclusive
to 180 exclusive. Zero omits the magnet. The magnet must not overlap a taper.
The arc is subdivided into slices with tangent magnetisation; the existing
|magnetization angle| is an additional local rotation of that direction.
Set it to 180 to reverse the magnet. This prescribed direction is transformed
with the node. Hysteresis history and irreversible demagnetisation are not added.
\end{key}
\begin{key}{/femm/component/magnet material=\meta{material} (initially cast-alnico-5-lng37)}
Catalogue or user material style for the toroidal magnet arc. Set
|ampere turns=0| for a magnet without a winding; otherwise both excitations
are present. The example below selects a linear-recoil NdFeB catalogue entry.
\end{key}
\begin{key}{/femm/component/arc step=\meta{positive number} (initially 5)}
Maximum angular span in degrees of polygon edges and magnet slices, at most
30, subject to a maximum of 4096 subdivisions per full revolution.
An excessively small value is rejected before allocating geometry.
Refine this independently of |mesh size| to improve the geometric and
magnetisation-direction approximation. It does not change a native TikZ path's
|curve tolerance|. Two same-material arc pieces can leave visible drawing seams.
\end{key}
\begin{stylekey}{/tikz/femm component magnet}
Base magnet painting style: |draw=red!70!black,fill=red!20|.
\end{stylekey}
\begin{key}{/femm/component/magnet style=\marg{TikZ options}}
Append local painting options to the magnet style, initially empty.
Use it as for the existing core, armature, coil and gap styles.
\end{key}

The toroids provide |core center| and |window center| at the annulus centre,
|upper pole| and |lower pole| at the two section centres, their four
|... pole inner|/|... pole outer| endpoints, |gap center| and |magnet center|.
The latter marks the opposite arc position even if no magnet is present.
A nonzero |ampere turns| adds one positive winding section in the window and
one negative return section outside the core, using the existing coil width,
height and clearance. The inner rectangle must fit strictly inside the window;
clearance must be positive. Zero omits both rectangles. The two coil anchors
exist only when these rectangles exist. |coil center y| is not used by toroids.
The component's envelope expands to include the exterior winding.

\luafemmexample{ideal-toroids}

\subsection{Discrete conservation, accuracy and interchange}
The confined domain is carved from an audited constrained mesh. Each outer
boundary has $A_z=0$ as a reference; all vertices of each window boundary
share one floating degree of freedom. The corresponding equation prescribes
the enclosed current, including explicit seeds and excluded winding regions.
Permanent magnets contribute through the usual constitutive weak form.
Independent material components have independent outer references.

Continuous P1 $A_z$ gives $\mathbf B=(\partial_y A_z,-\partial_x A_z)$.
Its normal component agrees across a shared triangle edge, and is zero on
an equipotential wall. Consequently any two consistently oriented complete
sections of one branch give the same endpoint potential difference, to
floating-point precision. Across a junction these differences telescope,
giving the signed flux balance. This remains true for nonlinear materials
and unequal sections. Local magnitudes, tangential continuity of $\mathbf H$
and contour circulations still need mesh-convergence checks.

Ideal field pictures space their levels over the boundary-potential range,
so the level spacing measures transported flux. A finite number of lines
does not encode every branch ratio as an exact ratio of integer line counts.
Permanent magnets can also produce local recirculation; confinement constrains
the net branch flux, not uniform pointwise induction. Explicit Lua contour
levels can inspect such additional potential extrema if needed.

Caches store the confined mesh and validate its boundary-potential ties on
restoration. Format 6 invalidates earlier private cache records. FEMM import
and FEM/ANS export are not offered for this mode: the floating-window/current
constraints are not represented by the existing interchange subset, and an
export must not silently turn the problem back into an open-air calculation.
The ordinary mode retains its existing interchange support.
